Análisis
El análisis de la relación entre el ingreso de los clientes y el saldo disponible en sus cuentas no evidencia un patrón lineal claramente definido. La dispersión de las observaciones muestra que clientes con niveles de ingreso similares pueden presentar saldos considerablemente diferentes y, de igual forma, clientes con ingresos elevados no necesariamente mantienen los saldos más altos. Este comportamiento se observa tanto en el conjunto de datos como al diferenciar los registros por segmento de clientes.
La distribución de los datos sugiere que el nivel de ingreso, por sí solo, no constituye un factor suficiente para explicar el comportamiento de los saldos bancarios. La ausencia de una tendencia claramente definida indica que podrían intervenir otros elementos que no se encuentran representados en la base de datos, tales como los hábitos de ahorro, la antigüedad del cliente, el uso de productos financieros, la actividad económica o características particulares de cada cliente.
A partir de este resultado, sería recomendable complementar el análisis mediante técnicas estadísticas, como el cálculo del coeficiente de correlación o la estimación de un modelo de regresión, con el fin de cuantificar la intensidad de la relación entre ambas variables y evaluar si esta resulta estadísticamente significativa. Asimismo, la incorporación de variables adicionales permitiría desarrollar un análisis multivariado que explique con mayor precisión los factores asociados al comportamiento de los saldos bancarios.